\begin{frame}
	\frametitle{Ответы на замечания Андреева~С.Н.}
	\textbf{Замечание 1 из 2}: В разделе 3.4 показано, что рассчитанные КФС воспроизводят асимметричную форму, наблюдаемую в эксперименте. Какой физический фактор в модели отвечает за асимметрию формы КФС, и подтверждается ли это независимыми экспериментальными данными?
	
	\hfill
	
	\textbf{Ответ:}
	
	
\end{frame}

\begin{frame}
	\frametitle{Ответы на замечания Андреева~С.Н.}
	\textbf{Замечание 1 из 2}: В главе 4 показано, что совместный учёт поправок на дрейф температуры окружающей среды и непостоянство коэффициента теплоотдачи снижает разброс между методиками с 25\% до 10\%. Из каких факторов складывается остаточная погрешность в 10\%?
	
	\hfill
	
	\textbf{Ответ:}
	
	
\end{frame}